% ============================================================
%  ANNOTATION TEMPLATE — Andre Ruiz Loera
%  Instructions:
%    1. Fill in the fields below.
%    2. Write your notes in the sections provided.
%    3. Compile with pdflatex or xelatex.
%    4. Also copy annotation-template.html, rename it to match
%       your book (e.g. rudin-principles.html), and paste your
%       content there so it shows on the website.
% ============================================================

\documentclass[11pt, letterpaper]{article}

\usepackage[margin=1.2in]{geometry}
\usepackage{amsmath, amssymb, amsthm}
\usepackage{mathtools}
\usepackage{hyperref}
\usepackage{enumitem}
\usepackage{parskip}
\usepackage{microtype}
\usepackage{booktabs}
\usepackage{xcolor}

% Theorem environments
\newtheorem{theorem}{Theorem}[section]
\newtheorem{lemma}[theorem]{Lemma}
\newtheorem{proposition}[theorem]{Proposition}
\newtheorem{corollary}[theorem]{Corollary}
\theoremstyle{definition}
\newtheorem{definition}[theorem]{Definition}
\newtheorem{example}[theorem]{Example}
\theoremstyle{remark}
\newtheorem*{remark}{Remark}
\newtheorem*{note}{Note}

% Convenience macros
\newcommand{\R}{\mathbb{R}}
\newcommand{\N}{\mathbb{N}}
\newcommand{\Z}{\mathbb{Z}}
\newcommand{\Q}{\mathbb{Q}}
\newcommand{\C}{\mathbb{C}}
\newcommand{\E}{\mathbb{E}}
\newcommand{\Prob}{\mathbb{P}}
\newcommand{\norm}[1]{\left\lVert#1\right\rVert}
\newcommand{\abs}[1]{\left|#1\right|}
\newcommand{\inner}[2]{\langle #1,\, #2 \rangle}

% ── FILL IN THESE FIELDS ────────────────────────────────────
\newcommand{\booktitle}{Title of the Book}
\newcommand{\bookauthor}{Author Name}
\newcommand{\bookyear}{Year Published}
\newcommand{\category}{Math | Finance | Philosophy | Books | Other}
\newcommand{\dateread}{Month Year}
% ────────────────────────────────────────────────────────────

\title{\booktitle \\ \large Annotations \& Notes}
\author{Andre Ruiz Loera \\ \small \textit{\category}}
\date{Read: \dateread}

\begin{document}
\maketitle
\tableofcontents
\newpage

% ── OVERVIEW ────────────────────────────────────────────────
\section{Overview}
% 2–4 sentences: what is this book about and why did you read it?


% ── KEY IDEAS ───────────────────────────────────────────────
\section{Key Ideas}
% Bullet the main concepts, theorems, or arguments.
\begin{itemize}
  \item
\end{itemize}

% ── CHAPTER NOTES ───────────────────────────────────────────
\section{Chapter Notes}

\subsection{Chapter 1 — }
% Notes, definitions, proofs, examples, questions.


% ── NOTABLE PASSAGES ────────────────────────────────────────
\section{Notable Passages}
% Direct quotes worth keeping.
\begin{quote}
  ``Paste quote here.'' \hfill --- \bookauthor
\end{quote}

% ── EXERCISES / PROBLEMS ────────────────────────────────────
\section{Exercises \& Problems}
% Work through selected exercises here.

\subsection*{Problem X.Y}
\textit{Statement of the problem.}

\begin{proof}
  % Your solution.
\end{proof}

% ── TAKEAWAYS ───────────────────────────────────────────────
\section{Takeaways}
% What will you remember in a year? What changed how you think?
\begin{itemize}
  \item
\end{itemize}

\end{document}
